% 
% Assignment 2 for Image Processing.
% Timothy van der Valk.
% 2336715.
%

% Preamble {{{1
\documentclass[10pt, a4paper]{article}
\usepackage{geometry}
\usepackage{graphicx}
\usepackage{array}
\usepackage{tabularx}
\usepackage{caption}
\usepackage{float}
\usepackage{xfrac}

\usepackage{pgfplots}
\pgfplotsset{compat=1.18}

\setlength{\parindent}{0pt}
\setlength{\tabcolsep}{2pt}

\newcommand{\imagecell}[2]{%
	\begin{minipage}[t]{\linewidth}
		\centering
		\includegraphics[
		height=45mm,
		keepaspectratio
		]{#1}\\
		{\scriptsize #2}
	\end{minipage}%
}

\newcommand{\bigimagecell}[3]{%
	\begin{minipage}[t]{\linewidth}
		\centering
		\includegraphics[
		height=#3,
		keepaspectratio
		]{#1}\\
		{\scriptsize #2}
	\end{minipage}%
}

\title{Image Processing\\Assignment 2}
\author{Timothy van der Valk\\2336715}

\begin{document}
\maketitle

\section*{Implementation} %{{{1

I have implemented a Binary and Grayscale Line Detection pipeline in the following way:

\begin{itemize}
	\item Smooth the image using a Gaussian kernel.
	\item Compute the Edge Magnitude image using horizontal and vertical Sobel kernels.
	\item Apply thresholding.
	\item Compute the \textbf{Normalized} Hough Transform using a fixed resolution of 512 pixels.
		I decided to use the normalized variant because the wideness of the tram image caused
		most vertical lines to not pass the peak threshold and be lost. My implementation also
		makes no distinction between the grayscale and binary implementation. Binary images will simply
		accumulate values of 255 for every pixel. Because peak-finding is relative to the max, this does not yield
		different results compared to incrementing by 1. Also, normalization evens out these differences in scale either way.
	\item Perform Hough Peak finding with a percentual peak threshold and Binary Closing.
	\item Perform Line Segment Detection In the line tracing code, I have decided to take the largest edge
		intensity in a neighborhood of 3x3 around the current pixel, to make it compatible with
		Canny Edge Detection which outputs only 1-pixel thin lines. This approach yields more stable
		lines for low-resolution images.
	\item Draw the line segments in \textbf{Red} on top of a darkened version of the original edge image (30\% brightness).
\end{itemize}

\section*{Task 1 | Binary Line Detection} %{{{1

See Figure \ref{fig:task1}. I have experimented with the peak threshold and the edge threshold.
The effect of the edge threshold is to include more edges, which can both strengthen existing edges, and
introduce many small and sporadic edges. A higher peak threshold will filter edges based on the strongest edges present.
The effect of this parameter is that it cuts away many off the sporadic edges, since they are very small and can not compete
with the long and dominant edges of the tram.

From the 3 parameter combinations, I would rate the middle one as the best balanced. It allows sufficient edges to
be included in the Hough Transform, but also cuts away more weaker edges due to a higher peak threshold. The result is a
strong fit on the dominant horziontal edges of the tram, and a decent fit on the doors and windows. The head of the tram is
harder to detect due to the slight curvature.

\section*{Task 2 | Grayscale Line Detection} %{{{1

See Figure \ref{fig:task2}. 
A big noticeable difference between the Binary and Grayscale pipelines is that
grayscale includes far more sporadic edges as a consequence of high-intensity and high-frequency regions of
blurred lines. In the first Grayscale image, you can see that many diagonal lines are rendered
in places with a high average magnitude (bright spot). This is both a strength and a weakness of
the Grayscale pipeline, because less edges tend to be lost due to thresholding.

In the Binary pipeline, many edges around the tram head are clipped due to the thresholding.
In the Grayscale pipeline, almost all three images contain one or more well fitted edges around 
the tram head.

In all three Grayscale edge images, I have kept all parameters except the blurring factor the same.
This is because the Grayscale pipeline is most susceptible to noise because it does not threshold away any
weak edges. In my opinion an optimal result is obtained for the image with gaussian kernel size 7 (middle).
This image shows both a strong fit on the key horizontal tram edges, and several more decent edges on 
tram head and doors. This makes it on par, or even better than the third image in the Binary pipeline.
It has a moderate blur that smoothens out high-frequency edges, while still retaining a high edge magnitude
that allows the Hough Transform to detect them, contrary to the third image with kernel size 11.
Minimum length was kept high similar to the Binary pipeline, this removed many of the smaller, sporadic edges.

\section*{Extra 1 | Hough Transform Counting }

As an extra I have used the code in my NormalizedHoughTransform() function to return the
maximum accumulator array for the input image. The result when applied to the tram image can be seen in Figure \ref{fig:hough_max}.
There is a clear hotspot in the center due to the width of the tram image. This is why I decided to use the Normalized Hough Transform in my report.

\begin{figure}[H]
	\centering
	\includegraphics[width=3.5cm]{results/maximum_tram.png}
	\caption{Hough Maximum Counts 512x512 for a full white image in the same dimensions as the tram image.}
	\label{fig:hough_max}
\end{figure}

\pagebreak

% kernel, sigma, threshold | peak%, length, gap, closing
% task1_good      : 3, 0.4, 75 | 0.5, 64, 4, 15
% task1_bad_gap   : 3, 0.4, 75 | 0.5, 32, 16, 5 (Fuses less lines together, allows shorter lines)
% task1_long      : 3, 0.4, 75 | 0.5, 128, 8, 7 (Longer lines, less fusion, mid gap)

\begin{figure}[H] % Clean images {{{2
	\begin{tabularx}{\textwidth}{X}
		\imagecell{results/tram\_contrast.png}{Reference Image.}
		\\[0.5cm]
		\imagecell{results/task1\_image1.png}{
			$\text{gaussianKernel}=3,\; \sigma=0.4,\;\text{peakThreshold}=50\%,\;\text{edgeThreshold}=128\\
			\text{minLen} = 64,\;\text{maxGap} = 4,\;\text{closeKernel} = 15$.
		}
		\\[0.5cm]
		\imagecell{results/task1\_image2.png}{
			$\text{gaussianKernel}=3,\; \sigma=0.4,\;\text{peakThreshold}=60\%,\;\text{edgeThreshold}=75\\
			\text{minLen} = 64,\;\text{maxGap} = 4,\;\text{closeKernel} = 15$.
		}
		\\[0.5cm]
		\imagecell{results/task1\_image3.png}{
			$\text{gaussianKernel}=3,\; \sigma=0.4,\;\text{peakThreshold}=45\%,\;\text{edgeThreshold}=100\\
			\text{minLen} = 64,\;\text{maxGap} = 4,\;\text{closeKernel} = 15$.
		}
	\end{tabularx}
	\captionof{table}{
	Binary Line Detection pipeline applied to contrast adjusted tram image.}
	\label{fig:task1}
\end{figure}

% kernel, sigma | peak%, min intensity, length, gap, closing
% task2_cross   : 3, 0.4 | 0.5, 50, 64, 4, 5
% task2_smooth7 : 7, 1.2 | 0.5, 50, 64, 4, 5
% task2_blur    : 11, 2  | 0.5, 40, 64, 4, 5

\begin{figure}[H] % Clean images {{{2
	\begin{tabularx}{\textwidth}{X}
		\imagecell{results/tram\_contrast.png}{Reference Image.}
		\\[0.5cm]
		\imagecell{results/task2\_cross.png}{$\text{gaussianKernel}=3,\; \sigma=0.4,\;\text{minIntensity}=50,\;\text{peakThreshold}=50\%\\\text{minLen} = 64,\;\text{maxGap} = 4,\;\text{closeKernel} = 5$.}
		\\[0.5cm]
		\imagecell{results/task2\_smooth7.png}{$\text{gaussianKernel}=7,\; \sigma=1.2,\;\text{minIntensity}=50,\;\text{peakThreshold}=50\%\\\text{minLen} = 64,\;\text{maxGap} = 4,\;\text{closeKernel} = 5$.}
		\\[0.5cm]
		\imagecell{results/task2\_blur.png}{$\text{gaussianKernel}=11,\; \sigma=2,\;\text{minIntensity}=50,\;\text{peakThreshold}=50\%\\\text{minLen} = 64,\;\text{maxGap} = 4,\;\text{closeKernel} = 5$.}
	\end{tabularx}
	\captionof{table}{
	Grayscale Line Detection pipeline applied to contrast adjusted tram image.}
	\label{fig:task2}
\end{figure}


\end{document}
